\thispagestyle{vazio}
\vspace*{8mm}
{\titlefontsemi\small\color{ebookPrim}COLEÇÃO ESCALA TRI}\par\vspace{4mm}
{\titlefontblack\fontsize{34}{36}\selectfont\color{ebookDark}Matemática\\para o ENEM}\par\vspace{5mm}
{\titlefont\Large\color{ebookPrim}Volume \numvolume\ — \titulovolume}\par\vspace{10mm}

\begin{tcolorbox}[enhanced,colback=ebookSoft,colframe=ebookLight,arc=6pt,boxrule=.6pt,left=10pt,right=10pt,top=8pt,bottom=8pt]
\small
\textbf{A coleção}\par
\begin{itemize}[leftmargin=1.3em,itemsep=1pt,topsep=3pt]
  \item \textbf{Volume 1} — Números, Contagem, Estatística e Probabilidade
  \item \textbf{Volume 2} — Geometria e Medidas
  \item \textbf{Volume 3} — Álgebra, Funções e Geometria Analítica
\end{itemize}
Juntos, os três volumes cobrem todos os objetos de conhecimento de Matemática e suas Tecnologias
previstos na Matriz de Referência do ENEM.
\end{tcolorbox}

\vfill
{\footnotesize\color{ebookMuted}
\textbf{Questões oficiais:} reproduzidas dos cadernos de prova do ENEM (caderno azul, 2.º dia, \anosbanco),
publicados pelo Instituto Nacional de Estudos e Pesquisas Educacionais Anísio Teixeira (INEP).
\textbf{Parâmetros da TRI e gabaritos:} Microdados do ENEM e gabaritos oficiais (INEP).
Detalhes em \hyperlink{creditos}{Fontes e créditos}.\par\medskip
\textbf{Questões inéditas, textos teóricos e resoluções:} elaborados especialmente para esta coleção,
com apoio de inteligência artificial e revisão matemática e pedagógica em várias etapas.\par\medskip
Edição 2026. Material de estudo independente, sem vínculo com o INEP ou com o Ministério da Educação.\par}
\clearpage
